% Architecture figure -- AAAI-27, spans both columns.
%
% NEEDS in the preamble, beyond the kit and booktabs/amsmath:
%     \usepackage{tikz}
%     \usetikzlibrary{arrows.meta,positioning,fit,backgrounds,calc}
% tikz is not on the kit's forbidden list, and it draws vectors rather than
% including an image file, so the kit's image-format and trim/clip rules do not apply.
%
% WHAT THIS DRAWS, and what it deliberately does not:
%   top-down  theta -> gate g(t) -> dendrite/membrane     (macro modulates micro)
%   bottom-up s(t-1) -> dendritic branch input            (spikes modulate the next branch step)
% There is NO arrow from the spiking layer into the Kuramoto layer, because the trained
% configuration has none: graph_feedback_strength = 0 and spike_pulse_gain = 0
% (hyperparameter.py:95,112), and the rollout raises if either is non-zero
% (sw240/rollout.py:13-15). The graph is built once from gamma and held for the whole
% rollout. Drawing a loop into the oscillators would show something the model does not do.
%
% Symbols match method.tex exactly: gamma, A, phi, theta, g(t), x(t), h, V, s(t), rho, Q.

\begin{figure*}[!t]
\centering
\begin{tikzpicture}[
  font=\small,
  box/.style   = {draw, rounded corners=2pt, align=center, inner sep=4pt, minimum height=8mm},
  stage/.style = {box, fill=black!4, text width=23mm},
  macro/.style = {box, fill=violet!10, draw=violet!60, text width=34mm},
  micro/.style = {box, fill=teal!10,   draw=teal!60,   text width=34mm},
  panel/.style = {box, fill=black!2, text width=31mm},
  flow/.style  = {-{Latex[length=2mm]}, thick},
  td/.style    = {-{Latex[length=2.6mm]}, very thick, violet!70},
  bu/.style    = {-{Latex[length=2.6mm]}, very thick, teal!70},
  lbl/.style   = {font=\scriptsize, align=center, inner sep=1pt}
]

% ---------------------------------------------------------------- left: encoding
\node[stage] (img)  {Image\\$128\times128$};
\node[stage, right=5mm of img] (enc) {Encoder\\8 filters};
\node[stage, right=5mm of enc] (gam) {Patch activations\\$\gamma$, $64\times64$};
\draw[flow] (img) -- (enc);
\draw[flow] (enc) -- (gam);

% ---------------------------------------------------------------- what gamma produces
\node[box, fill=black!4, above right=4mm and 7mm of gam, text width=30mm] (graph)
  {Coupling graph $A$\\[-1pt]{\scriptsize feature $+$ geodesic, top-512}};
\node[box, fill=black!4, below right=4mm and 7mm of gam, text width=30mm] (drive)
  {Sensory drive $\phi$\\[-1pt]{\scriptsize $\pi\tanh(W_\phi\bar{\boldsymbol u})$}};
\draw[flow] (gam.east) -- ++(3mm,0) |- (graph.west);
\draw[flow] (gam.east) -- ++(3mm,0) |- (drive.west);

% ---------------------------------------------------------------- centre: the two scales
\node[macro, right=10mm of graph] (kur)
  {\textbf{Macro-scale}\\[1pt]Kuramoto coordination\\[-1pt]
   {\scriptsize $\theta(t{+}1)=\theta+\Delta t(\omega+f\phi+\kappa\sin(\phi-\theta)+\tfrac{K}{N}\sum A\sin(\theta_j-\theta_i-\alpha))$}};
\node[micro, below=20mm of kur] (snn)
  {\textbf{Micro-scale}\\[1pt]Dendrites $\to$ membrane $\to$ spikes\\[-1pt]
   {\scriptsize $h\to V\to s(t)=g(t)\,\mathrm{ind}(V>\vartheta)$}};

\draw[flow] (graph.east) -- (kur.west);
\draw[flow] (drive.east) -| ([xshift=-6mm]kur.south west) -- ([xshift=-6mm]kur.south west) -- (kur.south west);

% top-down: gating
\draw[td] ([xshift=-9mm]kur.south) -- node[lbl, left=1pt, text=violet!80]
  {\textbf{Top-down}\\gating $g(t)$\\[-1pt]2-step delay\\$x=g\sin\theta$} ([xshift=-9mm]snn.north);

% bottom-up: spike feedback into the next branch input -- stays inside the micro scale,
% which is why it curves back into the same block and not up into the oscillators.
\draw[bu] (snn.east) .. controls +(12mm,6mm) and +(12mm,-6mm) ..
  node[lbl, right=1pt, text=teal!80] {\textbf{Bottom-up}\\spike feedback\\$s(t{-}1)\!\to\!$ branch} (snn.north east);

% ---------------------------------------------------------------- right: readout
\node[panel, right=22mm of kur] (sync)
  {\textbf{Phase synchrony}\\[-1pt]{\scriptsize $\theta$ clusters by object}};
\node[panel, right=22mm of snn] (corr)
  {\textbf{Spike correlation}\\[-1pt]{\scriptsize $\rho\!\to\!Q=\prod_a\max(\rho_a,0)$}};
\node[box, fill=black!4, below=7mm of corr, text width=31mm] (mask)
  {\textbf{Object masks}\\[-1pt]{\scriptsize connected components $\to$ greedy merge}};

\draw[flow, violet!70] (kur.east) -- (sync.west);
\draw[flow, teal!70] (snn.east) -- (corr.west);
\draw[flow] (corr) -- (mask);

% synchrony is a diagnostic, not the scored output
\node[lbl, below=1mm of sync, text=black!55] {diagnostic only};

\end{tikzpicture}
\caption{%
Signal flow. An image is encoded into patch activations $\gamma$, which fix both the coupling
graph $A$ and the sensory drive $\phi$ once per image. Oscillator phases $\theta$ evolve under
Kuramoto coupling on $A$ (macro-scale) and gate the spiking layer through $g(t)$, a two-step
delayed function of the mean patch phase (top-down). Within the spiking layer the previous spike
$s(t{-}1)$ re-enters the dendritic branch input (bottom-up). Object masks are read from
\emph{spike} correlations: $\rho$ over the settled window, combined across the four components
into the affinity $Q$, then connected components and a greedy merge. Phase synchrony is reported
as a diagnostic and is never the scored output.}
\label{fig:architecture}
\end{figure*}
